\section{离散分布与概率质量函数}
\label{sec:05-Probability/03-Random-Variables/02}

一家小面包店记了 30 天的到店人数：3 天没人，5 天 1 位，8 天 2 位，7 天 3 位，4 天 4 位，2 天 5 位，1 天 6 位。店主要决定明早备多少料，凭感觉说``通常两三位''没法折算成面粉的斤两。

把频率当成概率的估计，这份记录就变成一张表：0 对 \(0.1\)，1 对 \(0.167\)，2 对 \(0.267\)，如此往下。有了这张表，``明天不超过 2 位''是 \(0.1+0.167+0.267=0.534\)，备料方案可以按这个数往下算。上一节说分布是数轴上的一份概率分配；取值可数时，这份分配就是这样一张表。

\begin{definition}
离散随机变量 \(X\) 的概率质量函数（PMF）为
\[
p_X(x)=P(X=x),
\]
满足 \(p_X(x)\ge0\) 与 \(\sum_x p_X(x)=1\)，求和遍历 \(X\) 的全部可能取值。对任意 \(A\subseteq\R\)，\(P(X\in A)=\sum_{x\in A}p_X(x)\)。
\end{definition}

``质量''这个比方值得留着：总量为 1 的东西被分成若干份，摆在数轴的若干个点上。两条要求分别是非负与总量守恒，都直接来自概率公理，没有新内容。``离散''只要求取值集合可数，不要求是整数——后面会碰到取值为 \(1,1/2,1/3,\ldots\) 的例子。

\subsection{取值等间隔，不等于质量摊得均匀}

掷两枚公平骰子，令 \(S\) 为点数和。取值是 \(2,3,\ldots,12\)，共 11 个等距的整数。若以为每个各占 \(1/11\)，那就把``取值等距''和``质量均匀''混为一谈了。

36 种等可能的有序结果里，凑成 \(s\) 的有多少种？\(s=2\) 只有 \((1,1)\)；\(s=7\) 有六种；两端各只有一种。数完除以 36：
\[
p_S(s)=\frac{6-\abs{7-s}}{36},\qquad s=2,\ldots,12.
\]
质量在 \(s=7\) 处最厚（\(6/36\)），向两端递减到 \(1/36\)。\(7\) 这个数字本身没有魔力，它只是背后压着六种骰子组合，而 \(2\) 背后只压着一种。每个取值接到多少质量，取决于它的原像里装了多少基础结果——这正是上一节那张图的算术版本。

\subsection{同一个分布的两种画法}

累积分布函数换一个角度问：不问``恰好等于 \(x\)''，问``不超过 \(x\) 累计了多少''。

\begin{definition}
随机变量 \(X\) 的累积分布函数（CDF）为 \(F_X(x)=P(X\le x)\)。
\end{definition}

对离散变量它就是 PMF 的前缀和 \(F_X(x)=\sum_{t\le x}p_X(t)\)，图形上是一条右连续、单调不减、从 0 爬到 1 的阶梯。

\begin{center}
\begin{tikzpicture}[font=\footnotesize,>=stealth]
  \begin{scope}[x=0.44cm,y=0.4cm]
    \draw[->] (-1,0) -- (11.4,0) node[right] {\(s\)};
    \draw[->] (-1,0) -- (-1,7.0);
    \foreach \s/\c in {0/1,1/2,2/3,3/4,4/5,5/6,6/5,7/4,8/3,9/2,10/1}
      {\fill[blue!25] (\s-0.34,0) rectangle (\s+0.34,\c);
       \draw[blue!60!black] (\s-0.34,0) rectangle (\s+0.34,\c);}
    \node[above] at (5,6) {\(6/36\)};
    \node[below] at (0,-0.5) {\(2\)};
    \node[below] at (5,-0.5) {\(7\)};
    \node[below] at (10,-0.5) {\(12\)};
    \node at (5,8.4) {PMF：质量摆在点上};
  \end{scope}
  \begin{scope}[xshift=7.4cm,x=0.44cm,y=0.075cm]
    \draw[->] (-1.4,0) -- (11.4,0) node[right] {\(s\)};
    \draw[->] (-1.4,0) -- (-1.4,41);
    \draw[dashed,black!40] (-1.4,36) -- (11.2,36);
    \node[left] at (-1.5,36) {\(1\)};
    \draw[blue!70!black,thick]
      (-1.4,0) -- (0,0) -- (0,1) -- (1,1) -- (1,3) -- (2,3) -- (2,6) -- (3,6)
      -- (3,10) -- (4,10) -- (4,15) -- (5,15) -- (5,21) -- (6,21) -- (6,26)
      -- (7,26) -- (7,30) -- (8,30) -- (8,33) -- (9,33) -- (9,35) -- (10,35)
      -- (10,36) -- (11.2,36);
    \draw[<->,black!70] (5.35,15) -- (5.35,21);
    \node[right,black!70] at (5.5,18) {跳高 \(=p_S(7)\)};
    \node[below] at (0,-2.7) {\(2\)};
    \node[below] at (5,-2.7) {\(7\)};
    \node[below] at (10,-2.7) {\(12\)};
    \node at (5,45) {CDF：阶梯累计存量};
  \end{scope}
\end{tikzpicture}

\emph{左右两幅画的是同一个分布。左图第 \(s\) 根柱子的高度，等于右图阶梯在 \(s\) 处向上跳的高度——两种画法携带的信息完全一样，但右边那种对连续变量照样能用。}
\end{center}

跳跃高度这件事写成公式就是
\[
P(X=x)=F_X(x)-F_X(x^{-}),
\]
其中 \(F_X(x^{-})\) 是从左侧逼近的极限。区间概率则由差值给出，\(P(a<X\le b)=F_X(b)-F_X(a)\)。离散情形下端点马虎不得：\(P(5<S\le 8)=16/36\)，而 \(P(5\le S\le 8)\) 要把 \(P(S=5)=4/36\) 补回来，得 \(20/36\)。开闭之差正是那一级台阶的高度。

CDF 的价值不在于算得更快——离散情形下它反而绕。它的价值在于普适：有没有 PMF、有没有密度都不影响它有定义。下一节会看到连续变量的 CDF 是一条光滑上升的曲线，没有台阶；再往后还有既有台阶又有斜坡的混合分布。一套语言横跨三种世界，这是 PMF 做不到的。

\subsection{PMF 的形状由实验机制决定}

PMF 不是编出来的，它的样子反映实验怎么进行。\(n\) 次独立试验，每次成功概率 \(p\)，问恰好成功 \(k\) 次的概率。一条指定了成功位置的具体序列，概率是 \(p^k(1-p)^{n-k}\)；这样的序列有 \(\binom nk\) 条，彼此互斥。于是
\[
P(X=k)=\binom nk p^k(1-p)^{n-k},\qquad k=0,1,\ldots,n.
\]
组合数负责数序列，两个幂次负责给每条序列定价。\(n=3,p=1/2\) 代进去得到 \(1/8,3/8,3/8,1/8\)，与上一节直接枚举八个序列的结果对上了。这类``先辨认机制再写公式''的做法，\hyperref[ch:042]{``常见分布：不要背表，要先辨认随机机制''}会系统地做一遍。

\subsection{质量必须落在 0 与 1 之间}

\(p_X(x)\) 本身就是一个事件的概率，所以恒有 \(p_X(x)\le1\)。这一条看着平淡，却是离散与连续之间最容易翻船的地方：下一节的密度不受这个限制，它可以取 5、取 100，因为它压根不是某个事件的概率。做机器学习的人早晚会撞上这件事——连续变量的似然值大于 1、对数似然为正，都是合法的，不是代码写错了。

另一个常见误会是把``离散''等同于``取整数''。取值集合 \(\set{1,1/2,1/3,\ldots}\) 上放质量 \(p_X(1/n)=2^{-n}\)，总和为 1，是个正经的离散随机变量，取值既不等距也不是整数。反过来，身高四舍五入到厘米虽然取整数，建模时通常还是当作连续量的离散化更合适。判断标准不在于取了什么数值，而在于概率结构是从计数长出来的，还是为了方便硬分的箱。

到店人数、点数和、成功次数，这些量都能一个一个数过去。测量值不行：等车的时间、零件的直径、模型输出的分数，取值填满一整段区间，可数求和从根上就用不了。下一节换一套工具，把``摆在点上的质量''改成``抹在区间上的浓度''。
